\begin{figure}[htp!]
  \centering
  \includegraphics[scale=0.22]{Esercizio16.3/Esercizio.jpg} 
\end{figure}
Verifichiamo le energie di soglia del processo:
\begin{equation}
    \pi N\,\to\,\Lambda K
\end{equation}
L'energia nel centro di massa è:
\begin{equation}
    S=m_\pi^2+m_N^2+2m_NE_\pi\ge\left(\sum_im_i\right)^2 
\end{equation}
Quindi ottengo (formula generale per la prima tabella):
\begin{equation}
    E_\pi^{th}=\frac{1}{2m_N}\left[\left(\sum_im_i\right)^2-m_\pi^2-m_N^2\right]=0.9\,\mathrm{GeV}
\end{equation}
Togliendo la massa nel primo caso:
\begin{equation}
    K_\pi^{th}=0.75\,\mathrm{GeV}
\end{equation}

Per quanto riguarda la seconda tabella è analogo, infatti prendendo il terzo processo come esempio:
\begin{equation}
    NN\,\to\,\Lambda K N
\end{equation}
\begin{equation}
    S=2m_N^2+2m_NE_N\ge\left(\sum_im_i\right)^2 
\end{equation}
\begin{equation}
    E_N^{th}=\frac{1}{2m_N}\left[\left(\sum_im_i\right)^2-2m_N^2\right]=2.489\,\mathrm{GeV}
\end{equation}
\begin{equation}
    K_N=E_N-m_N=1.58\,\mathrm{GeV}
\end{equation}

Andando a vedere quali processi sono proibiti notiamo subito che:
\begin{equation}
    NN\,\to\,\Lambda\Lambda \qquad NN\,\to\,\Sigma\Sigma 
\end{equation}
Non sono permessi in quanto $\Lambda$ e $\Sigma$ hanno come stranezza $S=-1$ mentre i nucleoni $N$ hanno stranezza $s=0$. Quindi è presente un $|\Delta s|>1$ quindi non è possibile. In quanto la stranezza viene conservata nei processi forti ed elettromagnetici e nei processi deboli posso avere un $|\Delta s|=1$.